\documentclass[sigconf]{acmart}

\usepackage{graphicx}
\usepackage{booktabs}
\usepackage{url}

\settopmatter{printacmref=false}
\renewcommand\footnotetextcopyrightpermission[1]{}
\pagestyle{plain}

\title{Loom: An Open-Source Outfit Recommendation Engine\\with Semantic Scoring and Shopify Integration}

\author{Anushree Berlia}
\affiliation{%
  \institution{Independent Researcher}
}
\email{anushree.berlia@gmail.com}

\begin{document}

\begin{abstract}
We demonstrate \textbf{Loom}, an open-source outfit recommendation system
that generates complete, occasion-appropriate outfits from fashion catalogs.
Loom combines FashionCLIP embedding retrieval with a multi-signal scoring
function incorporating semantic material compatibility, vibe/anti-vibe
occasion priors, and style direction diversity.  The system is deployed in
two modes: (1)~a personal wardrobe app with weather-aware suggestions and
taste learning, and (2)~a Shopify storefront extension that adds ``Shop the
Look'' outfit recommendations to product pages.  Given a single anchor
garment, Loom generates three stylistically distinct outfits (Classic,
Trendy, Bold) in under 5 seconds on commodity hardware.  The system is
fully open-source and available at
\url{https://github.com/anushreeberlia/loom}.
\end{abstract}

\keywords{outfit recommendation, fashion AI, FashionCLIP, multimodal
embeddings, Shopify, recommendation systems}

\maketitle

\section{Introduction}

Outfit recommendation differs from single-item recommendation in that
the system must compose a \emph{set} of mutually compatible garments
satisfying visual, textual, and contextual constraints simultaneously.
Users expect outfits to respect color harmony, formality consistency,
occasion appropriateness, and material layering logic, constraints that
purely embedding-based systems struggle to enforce.

Prior systems address this through learned compatibility
embeddings~\cite{vasileva2018learning}, graph neural
networks~\cite{cui2019dressing}, or transformer-based set
completion~\cite{sarkar2023outfittransformer}.  However, these methods
capture statistical co-occurrence patterns but cannot enforce hard domain
constraints (e.g., a formal blazer should never pair with gym sneakers,
regardless of embedding proximity).  Industrial systems at Stitch Fix and
Zalando operate at scale but remain proprietary.

We present Loom, a fully open-source hybrid retrieval-scoring system that
addresses these challenges through:
\begin{itemize}
\item Slot-constrained ANN retrieval over FashionCLIP embeddings stored
  in PostgreSQL with pgvector (HNSW indexing).
\item A composite scoring function that integrates six signals: embedding
  similarity, color harmony, formality, occasion coherence, style direction,
  and diversity.
\item Two novel techniques: \emph{semantic material weight} (using CLIP
  geometry to infer garment heaviness) and \emph{vibe/anti-vibe occasion
  priors} (embedding prose occasion descriptions as scoring anchors).
\item Two deployment modes: a personal wardrobe app and a Shopify
  storefront extension for e-commerce merchants.
\end{itemize}

To our knowledge, Loom is the first open-source outfit recommendation
system that combines multimodal embedding retrieval with structured
fashion domain scoring and ships with both consumer and B2B deployment
targets.

\section{System Architecture}

\begin{figure}[t]
\centering
\fbox{\parbox{0.92\columnwidth}{\small\centering
\textbf{Loom Pipeline}\\[4pt]
\texttt{Anchor Image}
$\xrightarrow{\text{Fashion Florence}}$
\texttt{Structured Tags}
$\xrightarrow{\text{FashionCLIP}}$
\texttt{512-d Embedding}\\[3pt]
$\downarrow$\\[3pt]
$\forall\, s \in \{\text{bottom, shoes, layer, acc.}\}$:\\
\texttt{Query}$_s$ $\xrightarrow{\text{pgvector}}$ \texttt{Candidates}$_s$
$\xrightarrow{\text{filter}}$ \texttt{Scored}$_s$\\[3pt]
$\downarrow$\\[3pt]
\texttt{3 Outfits:} Classic $|$ Trendy $|$ Bold
}}
\caption{End-to-end pipeline from anchor image to three outfit outputs.}
\label{fig:pipeline}
\end{figure}

Loom processes an anchor item through four stages:

\textbf{(1) Vision.} Fashion Florence~\cite{berlia2026florence}, a
Florence-2-large model fine-tuned with LoRA, extracts structured attributes
(category, color, material, style tags, occasion tags) from the item image
as JSON.

\textbf{(2) Embedding.} FashionCLIP~\cite{chia2022fashionclip} produces a
512-dimensional blended embedding: $\mathbf{e} = \text{norm}(0.7 \cdot
\mathbf{e}_{\text{img}} + 0.3 \cdot \mathbf{e}_{\text{txt}})$, where
the text component encodes material, style, and category metadata.

\textbf{(3) Retrieval.} Per-slot cosine ANN queries retrieve candidates
from PostgreSQL via pgvector HNSW indexing ($m{=}16$,
$\text{ef\_construction}{=}64$).  Candidates are filtered by occasion
affinity using embedded vibe/anti-vibe prose vectors.

\textbf{(4) Scoring.} Combinatorial candidates are scored by a weighted
sum of embedding similarity, color harmony, formality consistency, occasion
coherence, and style direction bonuses.  The top outfit per direction
(Classic, Trendy, Bold) is selected.

\section{Key Technical Contributions}

\subsection{Semantic Material Weight}

Rather than maintaining a hand-coded material taxonomy, we compute garment
``heaviness'' in CLIP embedding space:
\begin{equation}
  w_i = \cos(\mathbf{e}_{\text{mat}(i)}, \mathbf{h}) -
        \cos(\mathbf{e}_{\text{mat}(i)}, \mathbf{l})
\end{equation}
where $\mathbf{h}$ and $\mathbf{l}$ are pre-embedded heavy/light context
strings.  A layer is compatible if $w_{\text{layer}} \geq w_{\text{top}}
- 0.15$.

\subsection{Vibe/Anti-Vibe Occasion Priors}

Each occasion (work, casual, going-out, etc.) is defined by two embedded
prose strings capturing what the occasion ``feels like'' and what it
explicitly rejects.  Items are scored by differential affinity:
\begin{equation}
  S_{\text{occ}}(i, o) = \cos(\mathbf{e}_i, \mathbf{v}_o) -
    \lambda_o \cdot \max(0, \cos(\mathbf{e}_i, \mathbf{a}_o) -
    \cos(\mathbf{e}_i, \mathbf{v}_o))
\end{equation}
New occasions require only writing new prose; no retraining is needed.

\section{Deployment Modes}

\subsection{Personal Wardrobe App}

Users photograph closet items, which are processed through the full
pipeline.  The system provides daily outfit suggestions that adapt to:
\begin{itemize}
\item \textbf{Weather}: OpenWeatherMap data adjusts material preferences
  (rain avoids suede, cold forces layers).
\item \textbf{Taste}: Like/dislike feedback updates an EMA-based taste
  vector that enters retrieval and scoring.
\item \textbf{Rotation}: A per-user FIFO queue penalizes recently worn
  items.
\end{itemize}
Authentication uses JWT with Google OAuth; data resides in PostgreSQL.

\subsection{Shopify Storefront Extension}

Merchants install Loom as an embedded Shopify app:
\begin{enumerate}
\item Catalog sync via Shopify Admin GraphQL API.
\item Automatic attribute extraction via Fashion Florence.
\item Pre-computed outfit generation with incremental cache invalidation.
\item ``Shop the Look'' theme extension on product detail pages.
\end{enumerate}

The Shopify integration uses React Router with Prisma session management.
Outfits are pre-generated and cached; adding a new product triggers
selective invalidation based on a category dependency map.

\section{Evaluation}

We evaluate Loom on an internal catalog of 620 items (230 tops, 170
bottoms, 140 shoes, 50 dresses, 30 accessories) sourced from H\&M
product data.

\subsection{Ablation Study}

\begin{table}[t]
\centering
\small
\caption{Ablation results (50 anchors, 150 outfits per config).}
\label{tab:ablation}
\begin{tabular}{lcc}
\toprule
Configuration & Score & Viol.\% \\
\midrule
Full system & 0.179 & 9.3 \\
$-$ Direction reranking & 0.052 & 6.7 \\
$-$ Material compat. & 0.233 & 4.7 \\
$-$ Occasion filtering & 0.194 & 7.3 \\
$-$ Formality check & 0.173 & 10.7 \\
Random retrieval & 0.054 & 16.0 \\
\bottomrule
\end{tabular}
\end{table}

Table~\ref{tab:ablation} summarizes ablation results.  The full system
achieves a 3.3$\times$ improvement over category-constrained random
retrieval (0.179 vs.\ 0.054) with 42\% fewer hard violations.  Direction
reranking is the single indispensable component: removing it drops the
score to 0.052, essentially equal to random.

\subsection{Latency}

\begin{table}[t]
\centering
\small
\caption{End-to-end latency breakdown (CPU, N=50 anchors).}
\label{tab:latency}
\begin{tabular}{lr}
\toprule
Stage & Time (ms) \\
\midrule
Vision (Fashion Florence, pre-computed) & --- \\
FashionCLIP embedding (ONNX, CPU) & 180 \\
pgvector retrieval ($5 \times 3$ queries) & 890 \\
Occasion filtering & 420 \\
Combinatorial scoring & 2,740 \\
\midrule
\textbf{Total (3 outfits)} & \textbf{4,395} \\
P95 & 5,354 \\
\bottomrule
\end{tabular}
\end{table}

Table~\ref{tab:latency} shows the latency breakdown.  The system
generates three complete outfits in 4.4 seconds on a single CPU core
without GPU acceleration.  Vision processing (Fashion Florence) is
pre-computed at catalog ingestion time.

\subsection{Diversity}

Across 50 anchor items, the system produces an average of 7.4 distinct
colors and 4.9 distinct category slots across the three style directions,
confirming meaningful aesthetic variation between Classic, Trendy, and Bold
outputs.

\section{Demonstration Scenarios}

The demonstration walks attendees through four scenarios that highlight
the system's capabilities:

\textbf{Scenario 1: Upload and Recommend.}
The user uploads a photo of a navy blazer.  Fashion Florence extracts
\texttt{\{category: "layer", color: "navy", material: "wool",
style\_tags: ["classic", "chic"], occasion\_tags: ["work"]\}}.  Loom
generates three outfits: a Classic look with tailored trousers and
oxford shoes, a Trendy option with dark jeans and white sneakers, and a
Bold combination with a patterned midi skirt and statement earrings.

\textbf{Scenario 2: Occasion Switching.}
Starting from the same navy blazer, the user switches the occasion from
``work'' to ``going-out.''  The system re-scores candidates using the
going-out vibe/anti-vibe vectors, producing outputs with sequined tops,
heeled boots, and clutch bags, items that would have been penalized under
the work occasion prior.

\textbf{Scenario 3: Taste Learning.}
The user ``likes'' outfits with earth tones and ``dislikes'' those with
bold prints.  The EMA taste vector shifts accordingly, and subsequent
recommendations favor muted palettes and solid fabrics.

\textbf{Scenario 4: Shopify ``Shop the Look.''}
A merchant's product detail page displays three complete outfit
recommendations below the main product image.  Each recommended item links
directly to its product page with an ``Add to Cart'' button, enabling
cross-sell conversion.

\medskip
\noindent A live demo is accessible via the Hugging Face Space at
\url{https://huggingface.co/spaces/anushreeberlia/fashion-florence}.

\noindent Source code: \url{https://github.com/anushreeberlia/loom}.

%% TODO: Add figure with screenshot grid showing the four scenarios
%% \begin{figure}[t]
%% \centering
%% \includegraphics[width=\columnwidth]{figures/demo_screenshots.pdf}
%% \caption{Screenshots of the four demonstration scenarios.}
%% \label{fig:screenshots}
%% \end{figure}

\section{Comparison with Existing Systems}

\begin{table}[t]
\centering
\small
\caption{Comparison with related outfit recommendation systems.}
\label{tab:comparison}
\begin{tabular}{lccccc}
\toprule
System & Open & Domain & Deploy & Occasion \\
       & Source & Rules & Modes & Aware \\
\midrule
Polyvore-based~\cite{vasileva2018learning} & \checkmark & -- & 0 & -- \\
OutfitTransformer~\cite{sarkar2023outfittransformer} & -- & -- & 0 & -- \\
DiFashion~\cite{xu2024difashion} & \checkmark & -- & 0 & -- \\
Stitch Fix & -- & \checkmark & 1 & \checkmark \\
\textbf{Loom (ours)} & \checkmark & \checkmark & 2 & \checkmark \\
\bottomrule
\end{tabular}
\end{table}

Table~\ref{tab:comparison} positions Loom among existing systems.  Loom
is unique in combining open-source availability, explicit domain
constraints, multiple deployment modes, and occasion-aware scoring.

\section{Conclusion}

Loom demonstrates that combining neural embedding retrieval with
structured domain scoring produces coherent, diverse outfits that respect
fashion constraints.  The system is fully open-source, runs on commodity
hardware, and is deployed in both personal and commercial (Shopify)
contexts.  All code, models, and evaluation scripts are publicly available.

\begin{acks}
This work uses FashionCLIP~\cite{chia2022fashionclip} and
Florence-2~\cite{xiao2024florence} as foundational components.
\end{acks}

\bibliographystyle{ACM-Reference-Format}
\begin{thebibliography}{7}

\bibitem[Berlia 2026]{berlia2026florence}
Anushree Berlia. 2026. Fashion Florence: Fine-tuning Florence-2 for
structured fashion attribute extraction. \emph{arXiv preprint}.

\bibitem[Chia et~al. 2022]{chia2022fashionclip}
Patrick~John Chia, Giuseppe Attanasio, Federico Bianchi, et~al. 2022.
FashionCLIP: Connecting language and images for product representations.
\emph{arXiv preprint arXiv:2204.03972}.

\bibitem[Cui et~al. 2019]{cui2019dressing}
Zeyu Cui, Zekun Li, Shu Wu, Xiaoyu Zhang, and Liang Wang. 2019.
Dressing as a whole: Outfit compatibility learning based on node-wise
graph neural networks. In \emph{WWW}.

\bibitem[Sarkar et~al. 2023]{sarkar2023outfittransformer}
Rohan Sarkar, Navaneeth Bodla, Mariya Vasileva, et~al. 2023.
OutfitTransformer: Learning outfit representations for fashion
recommendation. In \emph{WACV}.

\bibitem[Vasileva et~al. 2018]{vasileva2018learning}
Mariya~I. Vasileva, Bryan~A. Plummer, Krishna Dusber, et~al. 2018.
Learning type-aware embeddings for fashion compatibility. In \emph{ECCV}.

\bibitem[Xiao et~al. 2024]{xiao2024florence}
Bin Xiao, Haiping Wu, Weijian Xu, et~al. 2024. Florence-2: Advancing a
unified representation for a variety of vision tasks. In \emph{CVPR}.

\bibitem[Xu et~al. 2024]{xu2024difashion}
Yiyan Xu, Yongqi Yan, and Deng Lin. 2024.
DiFashion: Generative fashion outfit recommendation with diffusion
models. In \emph{SIGIR}.

\end{thebibliography}

\end{document}
